%==============================================================================
\section{Discussion}
\label{sec:discussion}

% --- maps to 01_Project_Overview.md §6 (contribution) and §7 (caveats) ---

\subsection{What information theory adds over $R^2$}
A correct, transferable noise metric: mutual information is valid where the
physics is non-linear and is expressed in bits. On the (PID) aging device the
measured dependence exceeds the linear-$R^2$ equivalent by $\valMIexcessPct\,\%$,
showing that $R^2$ systematically understates the true environment--sensor
coupling. This \emph{implies} that the prior $R^2$-based comparison of thermal
regimes understates their difference; recomputing both benches head-to-head in
bits would require a matched run on the same device, which is outside this
single-regime study (\cref{sec:discussion}, threats to validity).

\subsection{A principled definition of burn-in quality}
Reframing burn-in quality as the minimization of $\MI{\slk}{\Tdie,\Vcore}$ gives
the qualitative claim ``PID is a cleaner validation environment'' a rigorous unit
and an optimization target: every bit the environment writes into the slack
reading is a bit of aging information lost. We state this as a principle rather
than demonstrate it within one device; the closest empirical support here is the
distributional gap between the two benches (\cref{sec:res-kl}), measured on a
separate test device.

\subsection{A description-length view of the linear-vs-non-linear choice}
The MDL criterion of \cref{sec:th-mdl} that selects the linear PVT correction
(\cref{sec:res-correction}) is not an ad-hoc penalty on extra parameters: it is
a computable stand-in for the Kolmogorov-complexity intuition that the best
model is the \emph{shortest description} that still accounts for the data
(\cref{eq:kolmogorov}). Framed this way, the result is not merely ``the
non-linear correction is not worth it here''---it is that at this single,
narrow operating point (\valDutTempMean~$^\circ$C, \valDutVoltMean~V) the
data do not yet contain enough curvature information to justify paying for the
Arrhenius/super-linear model's extra parameters, even though the underlying
physics is genuinely non-linear (\cref{sec:background-r2}). This is consistent
with, not contradictory to, the non-linearity result of \cref{sec:res-mi}: MI
detects non-linear \emph{dependence} that $R^2$ misses, while MDL asks whether
a specific parametric non-linear \emph{form} pays for itself in \emph{this}
narrow-range dataset. A wider-range or multi-setpoint campaign, spanning enough
of the Arrhenius/BTI curvature to be compressible by the non-linear model, is
the natural condition under which MDL would flip its verdict.

\subsection{A resolution limit for SLM sensors}
The capacity and BSC results put concrete numbers on what the sensor can do:
$\valCapacityStates$ distinguishable aging states over the campaign, a daily
($\valBSCwindow$~h) aging alarm of $\Cbsc=\valBSCcap$~bits, and a $3\sigma$
detection floor of \valMinDetectRate~m-cnt/h reached after $\sim$\valMinObsTime~h.
These tell an SLM designer how finely the sensor resolves aging and how reliably
it can raise an alarm under realistic noise---quantities directly consumed by
adaptive voltage scaling and RUL, and which the prior residual-$\sigma$ metric
left implicit. Their mutual consistency (the BSC per-window reliability crosses
over near the CRLB minimum observation time) is evidence the channel model is
coherent rather than three disconnected heuristics.

\subsection{Sequential estimation and RUL as a distribution}
The Wiener/Kalman treatment of \cref{sec:res-wiener} is deliberately not a
replacement for the Cram\'er--Rao/Bayesian bound of \cref{sec:res-rul}, but a
cross-check from an independent estimation route (a dynamic, filtered model
rather than a single whole-series regression) that happens to agree on the
sign and order of magnitude of the aging rate. Its added value is
structural, not numerical: it turns a single point-rate-with-error-bar into
(i) a continuously updated, denoised trajectory estimate the operator could
in principle consume online, and (ii) a full RUL \emph{distribution} rather
than a bound, exposing the right-skew that a Gaussian error bar hides. That
the rate-noise MLE and the MDL criterion of \cref{sec:res-correction}
\emph{independently} both prefer the simplest available model (constant
drift; linear correction) at this single operating point is itself a
finding: two different estimation principles, applied to two different
questions, agree that this campaign does not yet contain the curvature
needed to justify a more flexible model --- exactly the kind of
cross-validation this project has used throughout (KSG vs.\ binned MI,
frequentist vs.\ Bayesian RUL, ICA vs.\ regression correction).

\subsection{Threats to validity}
The aging-device analysis is a single device and single thermal regime, which
limits permanent-aging separability. The bench comparison of \cref{sec:res-kl}
is on a \emph{different} test device, so its KL/JS divergence speaks to bench
fidelity, not to the aging device's behaviour, and the two are never mixed. MI
estimates on autocorrelated data are bias-prone despite the large raw sample
count, mitigated (not eliminated) by $\neff$, the surrogate null floor of
\cref{sec:res-mi} and cross-estimator agreement; the capacity expression is an
effective-resolution heuristic; the ICA separation is corroborative, not a clean
unmixing. The Wiener/Kalman model of \cref{sec:res-wiener} assumes constant
block-mean observation noise and a genuinely diffusive (Wiener) latent
process; the RUL distribution it produces is reported at self-referential,
already-computed thresholds precisely because no independently calibrated
failure threshold exists for this sensor, so it is a \emph{methodology}
demonstration, not an operational lifetime prediction, in the same spirit as
the Cram\'er--Rao bound of \cref{sec:res-rul}.
